\subsection{Real-Time Inter-Process Communication via UDP}\label{subsec:udp}

A real-time hand-tracking system that feeds a separate rendering
engine requires an inter-process communication mechanism capable of
sustaining a continuous, low-latency data stream.
This section discusses the choice of transport protocol for that link and
the use of a structured serialisation format.

\subsubsection{UDP vs.\ TCP for Real-Time Streaming}

The two dominant transport-layer protocols for socket communication
are the Transmission Control Protocol (\acrshort{tcp}) and the User
Datagram Protocol (\acrshort{udp}).

\acrshort{tcp} guarantees ordered, reliable delivery by
retransmitting any packet that is lost in transit.
In a continuous pose stream, however, this guarantee becomes
a liability: a retransmitted packet arrives late and carries data
that describes a hand position that has already changed.
Waiting for stale data introduces latency without adding information,
which is worse than simply skipping a frame.

\acrshort{udp} is connectionless and offers no delivery guarantees.
Each datagram is sent once and either arrives or is silently
discarded; the protocol never stalls to recover a lost packet.
The receiver always acts on the most recently arrived datagram,
so the worst consequence of packet loss is a single skipped frame
rather than a latency spike.
For applications where timeliness is more important than completeness
--- such as live pose streaming --- \acrshort{udp} is therefore the
appropriate choice.

\subsubsection{Structured Serialisation with JSON}

Each datagram must carry several heterogeneous values simultaneously:
joint rotation data for both hands, a scalar gesture label, and a
3-D anchor position.
Packing these into a single message requires a serialisation format
that is both compact enough for real-time use and flexible enough to
accommodate optional fields (e.g.\ when only one hand is visible).

\acrshort{json} satisfies both requirements.
It represents structured data as human-readable key-value pairs,
supports optional and nullable fields natively, and is natively
supported by virtually every programming language and environment.
Its text-based nature does carry a modest size overhead compared to
binary formats, but for the data volumes involved in a per-frame pose
message this overhead is negligible.
A datagram is only transmitted when at least one hand has been
detected, keeping the channel idle during frames with no visible
hands.
